\documentclass[11pt,a4paper]{article}

% ---- Moteur : compilé avec Tectonic (XeTeX) ----
\usepackage{fontspec}
\usepackage[french]{babel}
\usepackage{mathtools}
\usepackage{amssymb}
\usepackage[a4paper,margin=2.5cm,headheight=15pt]{geometry}
\usepackage{xcolor}
\usepackage{titlesec}
\usepackage{fancyhdr}
\usepackage{booktabs}
\usepackage{array}
\usepackage{tabularx}
\usepackage{enumitem}
\usepackage{tikz}
\usetikzlibrary{arrows.meta,positioning,calc}
\usepackage{caption}
\usepackage[most]{tcolorbox}
\usepackage[colorlinks=true,linkcolor=navy,urlcolor=goldd]{hyperref}

% ---- Couleurs de l'identité produit ----
\definecolor{navy}{RGB}{15,31,61}
\definecolor{navymid}{RGB}{26,46,82}
\definecolor{goldd}{RGB}{166,132,46}
\definecolor{goldl}{RGB}{201,164,76}
\definecolor{ink}{RGB}{26,34,48}
\definecolor{muted}{RGB}{90,100,120}
\definecolor{boxbg}{RGB}{244,247,251}
\definecolor{border}{RGB}{216,222,233}
\color{ink}

% ---- Titres ----
\titleformat{\section}{\Large\bfseries\color{navy}}{\thesection.}{0.6em}{}
\titleformat{\subsection}{\large\bfseries\color{goldd}}{\thesubsection}{0.6em}{}
\titlespacing*{\section}{0pt}{18pt}{8pt}
\titlespacing*{\subsection}{0pt}{12pt}{4pt}

% ---- En-tête / pied ----
\pagestyle{fancy}
\fancyhf{}
\renewcommand{\headrulewidth}{0.4pt}
\renewcommand{\footrulewidth}{0.4pt}
\fancyhead[L]{\small\bfseries\color{navy}LexIA Maroc}
\fancyhead[R]{\small\color{muted}Comprendre le projet}
\fancyfoot[C]{\small\color{muted}\thepage}

% Interligne un peu aéré pour la fluidité de lecture
\linespread{1.06}

% ---- Encadré « en clair » ----
\newtcolorbox{encadre}[1]{
  colback=boxbg, colframe=boxbg, boxrule=0pt, leftrule=2.6pt,
  borderline west={2.6pt}{0pt}{goldl},
  arc=3pt, left=12pt, right=12pt, top=9pt, bottom=9pt,
  fonttitle=\small\bfseries\color{navy}, coltitle=navy,
  title={#1}, attach title to upper=\par\vspace{3pt},
  fontupper=\small\itshape\color{ink}}

\begin{document}

% ═══════════════════ PAGE DE GARDE ═══════════════════
\begin{titlepage}
\thispagestyle{empty}
\centering
\vspace*{2.4cm}

\begin{tikzpicture}
  \node[rounded corners=8pt, fill=navy, minimum width=2cm, minimum height=2cm] (b) {};
  \node at (b) {\fontsize{40}{40}\selectfont\bfseries\color{goldl}\S};
\end{tikzpicture}

\vspace{1.1cm}
{\fontsize{40}{44}\selectfont\bfseries\color{navy} LexIA Maroc}\\[0.5cm]
{\Large\color{goldd} L'assistant juridique intelligent pour le droit marocain}\\[0.8cm]
{\color{goldl}\rule{0.4\textwidth}{1.2pt}}\\[0.8cm]
{\large\color{ink} Comprendre le projet}\\[0.25cm]
{\color{muted} Une explication simple de ce que fait la plateforme et de la façon dont elle fonctionne}\\[2.6cm]

\begin{tabularx}{0.82\textwidth}{@{}>{\bfseries\color{navy}}p{4cm} X@{}}
\toprule
Auteur & N'Guémawen N'DJINILA \\[3pt]
Cadre & Stage de fin d'année --- Transformation Digitale Industrielle \\[3pt]
Organisme & Zenithsoft, Rabat \\[3pt]
Version & LexIA Maroc v2.0 \\[3pt]
Contenu indexé & Plus de 6\,000 passages de textes juridiques officiels \\
\bottomrule
\end{tabularx}

\vfill
{\small\color{muted} Document composé avec \LaTeX.}
\end{titlepage}

% ═══════════════════ SOMMAIRE ═══════════════════
{\small\color{muted}\tableofcontents}
\clearpage

% ═══════════════════ 1. L'IDÉE ═══════════════════
\section{En deux mots : à quoi sert LexIA ?}

Imaginez pouvoir poser une question de droit marocain --- « Quelle est la durée de la période
d'essai d'un cadre ? », « Que risque un salarié licencié abusivement ? » --- et recevoir, en
quelques secondes, une réponse claire, rédigée en français, \textbf{qui cite précisément l'article
et la page du texte officiel} sur lequel elle s'appuie. C'est exactement ce que fait
\textbf{LexIA Maroc}.

Le droit marocain est vaste et éparpillé : Bulletin officiel, code du travail, code de commerce,
code des impôts, dahirs, décrets, circulaires… Retrouver la bonne disposition demande du temps et de
l'expertise. Les assistants d'intelligence artificielle grand public, eux, ne sont pas fiables sur ce
terrain : ils \emph{inventent} régulièrement des articles, ignorent les textes marocains et ne citent
aucune source vérifiable.

\begin{encadre}{L'idée centrale, en une phrase}
Plutôt que de faire confiance à la mémoire d'une intelligence artificielle --- qui se trompe parfois
avec aplomb ---, LexIA commence toujours par \textbf{retrouver les vrais textes officiels}, puis
demande à l'IA de rédiger une réponse \textbf{uniquement à partir de ces textes}, en citant ses
sources. L'IA ne « sait » rien : elle met en forme ce qu'on lui montre.
\end{encadre}

Cette approche porte un nom : le \textbf{RAG} (\textit{Retrieval-Augmented Generation}, c'est-à-dire
« génération assistée par la recherche »). C'est ce qui sépare un simple robot conversationnel d'un
assistant juridique digne de confiance : chaque affirmation est ancrée dans un texte réel, et
l'utilisateur peut la vérifier.

% ═══════════════════ 2. VUE D'ENSEMBLE ═══════════════════
\section{Comment ça marche, d'un coup d'œil}

Quand vous posez une question, elle traverse une petite chaîne de traitements. Chaque étape a un rôle
simple, et l'ensemble transforme une question en langage courant en une réponse sourcée.

\begin{center}
\begin{tikzpicture}[
  node distance=6mm and 7mm,
  box/.style={draw=border, fill=boxbg, rounded corners=3pt, minimum height=9mm,
              inner sep=4pt, font=\scriptsize\bfseries, text=navy, align=center, minimum width=20mm},
  arr/.style={-{Stealth[length=5pt]}, navymid, thick}]
  \node[box] (q) {Votre\\question};
  \node[box, right=of q] (r) {On choisit\\le domaine};
  \node[box, right=of r] (h) {On retrouve\\les passages};
  \node[box, right=of h] (rk) {On garde\\les meilleurs};
  \node[box, below=8mm of rk] (c) {On mesure\\la confiance};
  \node[box, left=of c] (g) {L'IA rédige\\la réponse};
  \node[box, left=of g] (s) {Réponse +\\sources};
  \draw[arr] (q)--(r); \draw[arr] (r)--(h); \draw[arr] (h)--(rk);
  \draw[arr] (rk)--(c); \draw[arr] (c)--(g); \draw[arr] (g)--(s);
\end{tikzpicture}
\captionof{figure}{Le parcours d'une question, de la saisie à la réponse sourcée.}
\end{center}

Les trois grandes phases sont : \textbf{préparer} les textes une fois pour toutes (en coulisses),
\textbf{retrouver} les passages utiles à chaque question, puis \textbf{rédiger} la réponse. Les
sections qui suivent racontent chacune de ces phases, sans jargon.

% ═══════════════════ 3. PRÉPARER LES TEXTES ═══════════════════
\section{En coulisses : préparer les textes}

Avant de pouvoir répondre, LexIA doit « digérer » les textes officiels. Un PDF de loi n'est pas
directement exploitable : il faut en extraire le texte, le découper en morceaux, et le ranger de
façon à pouvoir le retrouver vite.

\subsection{Lire les PDF, même les scans}
La plupart des textes sont lus directement (leur PDF contient déjà du texte). Certains, plus anciens,
sont de simples \emph{images} scannées : pour ceux-là, LexIA utilise la \textbf{reconnaissance
de caractères} (OCR), capable de lire aussi bien le français que l'arabe.

\subsection{Découper intelligemment}
Un texte juridique est naturellement organisé en articles. LexIA respecte cette structure et découpe
les textes \textbf{article par article}, plutôt que de trancher au milieu d'une phrase. Chaque
morceau garde en mémoire \textbf{sa page d'origine} --- c'est ce qui permet, plus tard, d'écrire
« \textit{Code du travail, article 62, page 34} ».

\subsection{Ranger dans deux « bibliothèques »}
Chaque morceau est rangé de deux manières complémentaires : une bibliothèque « par le sens » et une
bibliothèque « par les mots ». Pourquoi deux ? Parce qu'elles se rattrapent l'une l'autre, comme on
va le voir --- l'une comprend les idées, l'autre retrouve les mots exacts.

% ═══════════════════ 4. RETROUVER LES BONS PASSAGES ═══════════════════
\section{Le cœur du système : retrouver les bons passages}

C'est l'étape décisive. Une réponse ne vaut que par la qualité des passages sur lesquels elle
s'appuie. LexIA combine plusieurs techniques, chacune corrigeant les faiblesses des autres.

\subsection{Choisir le bon rayon}
Inutile de fouiller le code des impôts pour une question sur le licenciement. LexIA \textbf{devine
d'abord le domaine} de la question (travail, fiscal, sociétés, données personnelles) à partir des
mots employés, et concentre sa recherche là où les réponses ont le plus de chances de se trouver. En
cas de doute, il élargit la recherche à tous les domaines : mieux vaut chercher trop large que passer
à côté.

\subsection{Comprendre le sens : la recherche « par le sens »}
La première bibliothèque cherche par le \textbf{sens}, pas par les mots. Chaque texte reçoit une
sorte d'« empreinte de sens » : deux textes qui parlent de la même chose ont des empreintes proches,
\textbf{même s'ils n'emploient pas les mêmes mots}. Ainsi, « licenciement abusif » et « rupture
injustifiée du contrat » sont reconnus comme voisins, là où une simple recherche de mots-clés
échouerait. C'est ce qui donne à LexIA sa capacité à « comprendre » une question formulée avec vos
propres mots.

\subsection{Retrouver les mots exacts : la recherche « par les mots »}
Mais le sens ne suffit pas toujours. Si vous cherchez « article 62 » ou un numéro de loi précis, il
faut une correspondance \textbf{à la lettre}. La seconde bibliothèque s'en charge : c'est une
recherche par mots-clés, à la manière d'un moteur de recherche classique, avec une finesse en plus
--- elle sait qu'un \textbf{mot rare} (un terme juridique pointu) est bien plus révélateur qu'un mot
courant, et lui donne davantage de poids.

\begin{encadre}{Un détail qui compte}
Certains mots-clés très courts, comme « IS » (impôt sur les sociétés) ou « IR » (impôt sur le
revenu), se cachent à l'intérieur de mots français ordinaires --- « préav\textbf{is} »,
« secréta\textbf{ir}e ». Sans précaution, ils déclenchaient de mauvais aiguillages. LexIA ne les
reconnaît donc que lorsqu'ils forment un mot entier.
\end{encadre}

\subsection{Faire voter les méthodes}
LexIA lance ces recherches en parallèle (et même sur plusieurs reformulations de la question, pour ne
rien manquer). Il obtient donc plusieurs listes de résultats. Comment les combiner ? Par un
\textbf{vote} : un passage bien classé par \emph{plusieurs} méthodes à la fois remonte en tête. Cette
façon de fusionner est robuste, car elle ne se fie pas aux « notes » de chaque méthode (difficiles à
comparer entre elles), mais seulement à l'\textbf{ordre} dans lequel elles classent les passages.

\subsection{Un second lecteur, plus fin}
Les premières méthodes sont rapides mais un peu grossières : elles jugent la question et chaque
passage \emph{séparément}. LexIA fait alors intervenir un \textbf{second lecteur}, plus lent mais
beaucoup plus fin, qui relit la question \textbf{et} le passage \textbf{ensemble} pour juger leur
pertinence réelle. Comme il est coûteux, on ne lui confie que la petite dizaine de meilleurs
candidats. Un soin particulier est apporté aux \textbf{numéros d'article} : si vous demandez
l'article 52 et qu'un passage \emph{est} l'article 52, il est mis en avant.

\subsection{Savoir dire « je ne suis pas sûr »}
À l'issue de ce tri, LexIA attribue à la réponse un \textbf{niveau de confiance} (élevé, moyen,
faible). C'est une qualité rare et précieuse : si les meilleurs passages trouvés sont trop faiblement
liés à la question, le système \textbf{préfère l'admettre} --- « aucun texte pertinent trouvé,
reformulez ou précisez le domaine » --- plutôt que d'inventer une réponse hasardeuse.

% ═══════════════════ 5. RÉDIGER LA RÉPONSE ═══════════════════
\section{Rédiger la réponse}

Les meilleurs passages sont enfin transmis à un modèle de langage (le « cerveau rédacteur »). Son
rôle est volontairement limité : il ne puise pas dans sa mémoire, il \textbf{reformule et structure}
les extraits qu'on lui fournit.

\subsection{Une réponse adaptée à qui la lit}
La même question n'appelle pas la même réponse selon qu'elle vient d'un étudiant ou d'un avocat.
LexIA adapte donc le ton et la structure au \textbf{profil} de l'utilisateur.

\begin{table}[h]
\centering\small
\begin{tabularx}{\textwidth}{@{}>{\bfseries\color{navy}}p{3.2cm} X@{}}
\toprule
Étudiant & Pédagogique : on définit les termes, on explique le principe. \\
Particulier & Accessible et concret, orienté « que dois-je faire ? ». \\
Juriste / Avocat & Technique et précis, avec les références d'articles. \\
Entreprise & Orienté décision : obligations, risques, actions à mener. \\
\bottomrule
\end{tabularx}
\end{table}

\subsection{Des garde-fous contre l'invention}
Plusieurs règles strictes encadrent la rédaction : LexIA traite \textbf{exclusivement} le droit
marocain (il ne doit jamais présenter un texte comme étranger --- il lui est arrivé, avant correction,
de qualifier une loi marocaine « de loi française ») ; il ne doit rien affirmer qui ne figure pas dans
les passages fournis ; et lorsqu'il ne trouve pas, il indique ce que couvrent les textes et invite à
reformuler, au lieu d'un « aucune information » sec. C'est cette discipline qui fait la fiabilité de
l'ensemble.

% ═══════════════════ 6. LES OUTILS ═══════════════════
\section{Bien plus qu'un chat : les outils}

Autour de la conversation, LexIA propose des outils concrets qui en font une véritable plateforme.

\begin{itemize}[leftmargin=1.3em,itemsep=4pt]
  \item \textbf{Calculateurs} --- indemnité de licenciement, préavis, salaire net. Ces calculs sont
        exacts et reproductibles (ils appliquent le barème légal, sans intervention de l'IA) et citent
        l'article applicable.
  \item \textbf{Audit de contrat} --- on importe un contrat, LexIA le confronte au code du travail et
        signale les points de conformité, les risques et les clauses manquantes. Il examine le contrat
        \textbf{thème par thème} (période d'essai, préavis, salaire, congés…) pour ne rien laisser
        passer.
  \item \textbf{Génération de documents} --- attestation de travail, mise en demeure, contrat CDI,
        reçu pour solde de tout compte : on remplit un formulaire, on télécharge le document en Word ou
        en PDF.
  \item \textbf{Comparaison de versions} --- visualiser ce qui a changé entre deux versions d'un même
        texte (par exemple d'une année de loi de finances à l'autre).
  \item \textbf{Veille réglementaire} --- LexIA surveille l'arrivée de nouveaux textes et les signale ;
        un bandeau « derniers textes intégrés » anime la page d'accueil.
  \item \textbf{Français et arabe} --- questions et réponses dans les deux langues, avec une interface
        adaptée à l'écriture de droite à gauche.
\end{itemize}

% ═══════════════════ 7. CONFIANCE ═══════════════════
\section{Sécurité et confiance}

Un service qui manipule des comptes utilisateurs doit être irréprochable. Les mots de passe ne sont
jamais stockés en clair ; les sessions reposent sur des jetons sécurisés qui ne peuvent pas être
volés facilement par une page malveillante ; l'accès est limité en fréquence pour prévenir les abus ;
et, en production, tout transite par une connexion chiffrée (HTTPS). Enfin, chaque réponse rappelle
qu'elle est \textbf{informative} et ne remplace pas l'avis d'un professionnel du droit.

% ═══════════════════ 8. UN VRAI PRODUIT ═══════════════════
\section{Un vrai produit, pas une maquette}

LexIA a été pensé pour vivre en conditions réelles. Il propose un \textbf{modèle d'abonnement} :

\begin{table}[h]
\centering\small
\begin{tabularx}{\textwidth}{@{}>{\bfseries\color{navy}}p{2.4cm} p{1.8cm} X@{}}
\toprule
Offre & Prix & Ce qu'elle inclut \\
\midrule
Gratuit & 0 MAD & 20 questions par mois, chat sourcé, calculateurs, recherche par référence. \\
Pro & 99 MAD & Questions illimitées, audit de contrat, génération de documents, export. \\
\bottomrule
\end{tabularx}
\end{table}

Toute la mécanique (quotas, limitation des fonctionnalités selon l'offre, compteur d'usage) est en
place ; il ne reste qu'à brancher un moyen de paiement réel. Côté technique, l'application est
entièrement \textbf{conteneurisée} (prête à être déployée d'un seul geste sur un serveur) et couverte
par une large batterie de tests automatiques. Plusieurs anomalies n'ont d'ailleurs été détectées
qu'en \textbf{utilisant réellement l'application} --- et non par les seuls tests --- puis corrigées :
c'est le signe d'un produit éprouvé, pas seulement codé.

% ═══════════════════ 9. BILAN ═══════════════════
\section{Bilan et prochaines étapes}

LexIA Maroc tient sa promesse : des réponses juridiques \textbf{fiables, sourcées et vérifiables},
adaptées au profil de chacun, complétées par des outils utiles au quotidien. Il reste des chantiers
pour aller plus loin :

\begin{itemize}[leftmargin=1.3em,itemsep=3pt]
  \item Brancher le \textbf{paiement en ligne} (la mécanique d'abonnement est prête).
  \item Ajouter les \textbf{conditions générales} et la politique de confidentialité.
  \item Enrichir le corpus avec la \textbf{jurisprudence} (les décisions de justice), encore absente.
\end{itemize}

% ═══════════════════ ANNEXE ═══════════════════
\clearpage
\section*{\color{navy}Annexe --- Pour les curieux : les formules}
\addcontentsline{toc}{section}{Annexe --- Les formules}
{\small\color{muted} Cette annexe s'adresse au lecteur technique. Elle n'est pas nécessaire pour
comprendre le reste du document.}
\vspace{6pt}

\noindent\textbf{\color{navy}Proximité de sens (similarité cosinus).} La ressemblance entre une
question \(\mathbf{q}\) et un passage \(\mathbf{d}\), une fois transformés en vecteurs, se mesure par
l'angle entre eux :
\[
\operatorname{sim}(\mathbf{q},\mathbf{d}) = \cos(\theta)
= \frac{\mathbf{q}\cdot\mathbf{d}}{\lVert\mathbf{q}\rVert\,\lVert\mathbf{d}\rVert}.
\]

\noindent\textbf{\color{navy}Recherche par les mots (BM25).} Le score de pertinence lexicale d'un
passage \(D\) pour une requête \(Q\), qui valorise les mots rares et corrige la longueur du passage :
\[
\operatorname{score}(D,Q) = \sum_{i} \operatorname{IDF}(q_i)\cdot
\frac{f(q_i,D)\,(k_1+1)}{f(q_i,D) + k_1\!\left(1 - b + b\,\frac{\lvert D\rvert}{\text{avgdl}}\right)}.
\]

\noindent\textbf{\color{navy}Le vote (fusion par rangs réciproques).} On additionne, pour chaque
passage, l'inverse de son rang dans chaque liste ; un passage bien classé partout gagne :
\[
\operatorname{RRF}(d) = \sum_{r} \frac{1}{k + \operatorname{rang}_r(d)}, \qquad k = 60.
\]

\noindent\textbf{\color{navy}Le niveau de confiance.} Le score du meilleur passage est ramené entre
0 et 1, puis traduit en label (élevé au-dessus de 0,8 ; moyen ; faible ; insuffisant en dessous de
0,4) :
\[
\sigma(s) = \frac{1}{1 + e^{-s}}.
\]

\vspace{12pt}
\noindent\textcolor{border}{\rule{\textwidth}{0.5pt}}\\[4pt]
{\small\color{muted} Document du projet LexIA Maroc --- Stage de fin d'année, Zenithsoft (Rabat). Les
réponses fournies par la plateforme sont informatives et ne constituent pas un avis juridique
professionnel.}

\end{document}
